\begin{proof}[Proof of \cref{obj:thm:marked-component-certificate}]
Fix the public parameters. All constants constructed below depend only on these parameters. We use the functions and sign arrays prescribed by \cref{obj:def:marked-handle}, extending the displayed function formulas to \(\mathbb R^d\) when taking analytic bounds.

\begin{enumerate}
\item \textit{Regularity and membership.}
Set
\[
D_d:=2^d.
\]
The function \(\chi\) is smooth and flat at zero: each of its derivatives on the positive half-line is a polynomial in \(1/u\) multiplied by \(\exp(-1/u)\), and this expression tends to zero as \(u\downarrow0\). The denominator defining the scalar angular transition in \cref{obj:def:marked-handle} is positive everywhere. Consequently the prescribed sine and cosine pieces join smoothly, including at their support endpoints. The bump \(B\) is also smooth and compactly supported, and
\[
0\le B_h(x)\le1,\qquad B_h(x_0)=1.
\]
Define the finite derivative bounds
\[
K_\omega:=1+\sup_{u\in\mathbb R}|\omega_0'(u)|,
\qquad
K_B:=1+\sup_{w\in\mathbb R^d}\|\nabla B(w)\|_2,
\]
and put
\[
K_\phi:=K_B+dK_\omega,
\qquad
K_F:=2D_dK_\phi.
\]
Translation invariance gives the same derivative bound for every \(\omega_z\). Telescoping a product of factors bounded by one, and then using \(\delta\le h\), yields
\[
|\phi_{\mathbf z}(x)-\phi_{\mathbf z}(x')|
\le
\left(\frac{K_B}{h}+\frac{dK_\omega}{\delta}\right)
\|x-x'\|_2
\le
\frac{K_\phi}{\delta}\|x-x'\|_2.
\]

At each point, a nonzero frame entry has, in each coordinate, one of the two lattice indices adjacent to the scaled coordinate. Thus at most \(D_d\) entries are nonzero. Each entry has absolute value at most one, so
\[
|F_\lambda(x)|\le D_d,
\qquad
|G_\eta(x)|\le D_d.
\]
For a difference between two points, entries outside the union of their active index sets contribute zero. That union has at most \(2D_d\) elements. Therefore
\[
|F_\lambda(x)-F_\lambda(x')|
\le \frac{K_F}{\delta}\|x-x'\|_2,
\qquad
|G_\eta(x)-G_\eta(x')|
\le \frac{K_F}{\delta}\|x-x'\|_2.
\]
For a Hölder order \(s_{\mathrm{ord}}\in(0,1]\), define
\[
C_{F,s_{\mathrm{ord}}}:=2D_d+K_F+1.
\]
When \(\|x-x'\|_2\le\delta\), the Lipschitz bound and
\[
\frac{\|x-x'\|_2}{\delta}
\le
\left(\frac{\|x-x'\|_2}{\delta}\right)^{s_{\mathrm{ord}}}
\]
give the required Hölder difference bound. When \(\|x-x'\|_2>\delta\), use the bound \(2D_d\) instead. Since \(\delta\le1\), the supremum term is covered as well. With the norm of \cref{obj:def:holder-norm}, this proves
\[
\|F_\lambda\|_{H^{s_{\mathrm{ord}}}}
\le C_{F,s_{\mathrm{ord}}}\delta^{-s_{\mathrm{ord}}},
\qquad
\|G_\eta\|_{H^{s_{\mathrm{ord}}}}
\le C_{F,s_{\mathrm{ord}}}\delta^{-s_{\mathrm{ord}}}.
\]

For the macro bump, use its fixed compact profile \(B^2\). For every integer \(j\ge0\), define
\[
D_{B,j}:=
1+\sup_{w\in\mathbb R^d}\|\mathrm D^j(B^2)(w)\|_{\mathrm{op}},
\]
where \(\mathrm D^j\) is the total derivative of order \(j\), with the order-zero norm interpreted as absolute value. These constants are finite by smoothness and compact support. For \(s_{\mathrm{ord}}>0\), define
\[
j_{\mathrm{der}}:=\lceil s_{\mathrm{ord}}\rceil-1,
\qquad
\xi_{\mathrm{der}}:=s_{\mathrm{ord}}-j_{\mathrm{der}}\in(0,1],
\]
\[
M_{B,s_{\mathrm{ord}}}
:=
2\max\{D_{B,j_{\mathrm{der}}},D_{B,j_{\mathrm{der}}+1}\},
\]
\[
C_{B,s_{\mathrm{ord}}}
:=
\sum_{j=0}^{j_{\mathrm{der}}}D_{B,j}d^j
+
M_{B,s_{\mathrm{ord}}}d^{j_{\mathrm{der}}}+1.
\]
The derivative of order \(j_{\mathrm{der}}\) has both a bounded supremum and a Lipschitz bound supplied by the next derivative. Splitting at distance one gives
\[
\|\mathrm D^{j_{\mathrm{der}}}(B^2)(w)
-\mathrm D^{j_{\mathrm{der}}}(B^2)(w')\|_{\mathrm{op}}
\le
M_{B,s_{\mathrm{ord}}}\|w-w'\|_2^{\xi_{\mathrm{der}}}.
\]
Dilation by \(h\) multiplies a derivative of order \(j\) by \(h^{-j}\). Evaluating a derivative on coordinate directions bounds a partial derivative by its operator norm times \(d^j\). Hence, for \(0<h\le1\),
\[
|\partial^\kappa(B_h^2)(x)|
\le D_{B,|\kappa|}d^{|\kappa|}h^{-|\kappa|}
\le C_{B,s_{\mathrm{ord}}}h^{-s_{\mathrm{ord}}}
\qquad (|\kappa|\le j_{\mathrm{der}}),
\]
and
\[
|\partial^\kappa(B_h^2)(x)-\partial^\kappa(B_h^2)(x')|
\le
M_{B,s_{\mathrm{ord}}}d^{j_{\mathrm{der}}}
h^{-s_{\mathrm{ord}}}\|x-x'\|_2^{\xi_{\mathrm{der}}}
\qquad (|\kappa|=j_{\mathrm{der}}).
\]
Thus, in particular for the orders \(\beta\) and \(\gamma\),
\[
\|B_h^2\|_{H^{s_{\mathrm{ord}}}}
\le C_{B,s_{\mathrm{ord}}}h^{-s_{\mathrm{ord}}}.
\]

Define
\[
K_M:=
D_d+C_{F,\alpha}+C_{F,\beta}
+C_{B,\gamma}+C_{B,\beta}+1,
\qquad
\rho:=\min\left\{1,L-1/2,\frac12-\varepsilon\right\},
\qquad
c_M:=\frac{\rho}{32K_M}.
\]
The public domain gives \(0<\rho\le1\). Hence \(0<c_M\le\rho/32\le1/32\), and multiplying any of the five summands preceding the final \(1\) in \(K_M\) by \(c_M\) gives at most \(\rho/32\).

Suppose the amplitude and contrast inequalities hold with \(c_M\). Since \(0<\delta\le h\le1/2\),
\[
a\le c_M,\qquad b\le c_M,\qquad ab\le c_M,
\]
\[
a\delta^{-\alpha}\le c_M,\qquad
b\delta^{-\beta}\le c_M,\qquad
ab h^{-\gamma}\le c_M.
\]
Also \(\beta\le\gamma\) and \(h\le1\), so
\[
ab\le c_Mh^\gamma\le c_Mh^\beta,
\qquad
ab h^{-\beta}\le c_M.
\]
The prescribed probabilities therefore satisfy
\[
|e_\lambda(x)-1/2|
\le aD_d\le\frac{\rho}{32}\le\frac14,
\qquad
aD_d\le\frac{\rho}{32}\le\frac12-\varepsilon,
\]
\[
|\mu_{0,\theta,\eta}(x)-1/2|,
\ |\mu_{1,\theta,\eta}(x)-1/2|
\le bD_d+ab
\le\frac{\rho}{32}+c_M\le\frac1{16}\le\frac38.
\]
These bounds imply the overlap and both probability-valued arm-mean conditions.

Subadditivity and homogeneity of the Hölder norm give
\[
\|aF_\lambda\|_{H^\alpha}\le C_{F,\alpha}c_M,
\qquad
\|bG_\eta\|_{H^\beta}\le C_{F,\beta}c_M,
\]
\[
\|\tau_\theta\|_{H^\gamma}
\le2C_{B,\gamma}c_M,
\qquad
\|\theta abB_h^2\|_{H^\beta}
\le C_{B,\beta}c_M.
\]
Consequently
\[
\|e_\lambda\|_{H^\alpha}
\le\frac12+\frac{\rho}{32}\le L,
\]
\[
\|\mu_{0,\theta,\eta}\|_{H^\beta}
\le\frac12+\frac{\rho}{16}\le L,
\qquad
\|\tau_\theta\|_{H^\gamma}\le\frac{\rho}{16}\le L.
\]
All the prescribed functions are Borel. The independent conditional Bernoulli construction gives uniform covariates and conditional exchangeability, and its conditional margins are exactly the prescribed functions. Together with the preceding bounds, these facts establish every condition in \cref{obj:def:primitive-class}.

To identify the canonical contrast, the prescribed control mean is continuous, and the prescribed treated mean is continuous because it is the sum of the control mean and \(\tau_\theta\). Their admissible designated versions agree with them almost surely by uniqueness of conditional margins. Two continuous versions that agree almost everywhere under the uniform design agree everywhere on the cube: a nonzero continuous difference would persist on a relative neighborhood of positive uniform measure, including at a boundary point. Thus
\[
P_{\theta,\lambda,\eta}\in\mathcal M,
\qquad
\tau_{P_{\theta,\lambda,\eta}}(x)
=
\mu_{1,\theta,\eta}(x)-\mu_{0,\theta,\eta}(x)
=
\tau_\theta(x)
\quad(x\in[0,1]^d).
\]
% lean: CausalSmith.Stat.TwosamplePointcateAnnotationFrontier.marked_membership

\item \textit{Normalization and singleton cancellation.}
Define expectation over the finite sign arrays by
\[
\mathbb E_\theta Q
:=
\sum_{\lambda,\eta}
w_\theta(\lambda,\eta)Q(\lambda,\eta).
\]
At each index, direct summation over the four sign pairs gives
\[
\frac{1+\theta l s_\eta/2}{4}\ge0
\quad(l,s_\eta\in\{-1,+1\}),
\qquad
\sum_{l,s_\eta\in\{-1,+1\}}\frac{1+\theta l s_\eta/2}{4}=1.
\]
Multiplying these identities over the finite index set proves
\[
w_\theta(\lambda,\eta)\ge0,
\qquad
\sum_{\lambda,\eta}w_\theta(\lambda,\eta)=1.
\]
The same finite-product summation gives
\[
\mathbb E_\theta\lambda_{\mathbf z}
=
\mathbb E_\theta\eta_{\mathbf z}=0,
\qquad
\mathbb E_\theta[\lambda_{\mathbf z}\eta_{\mathbf w}]
=
\begin{cases}
\theta/2,&\mathbf z=\mathbf w,\\
0,&\mathbf z\ne\mathbf w.
\end{cases}
\]
Indeed, matching indices use the scalar sign-pair correlation, while distinct indices include a zero-mean scalar factor.

The sine-cosine identity on each lattice interval gives a sum of squares equal to one. Tensoring it and multiplying by \(B_h^2\) yields
\[
\sum_{\mathbf z\in\mathcal I}\phi_{\mathbf z}(x)^2=B_h(x)^2.
\]
For \(x\in C_h\), all nonzero tensor entries lie in the prescribed finite index set; for \(x\notin C_h\), both sides vanish. Expanding the finite signed sums therefore gives
\[
\mathbb E_\theta F_\lambda(x)
=
\mathbb E_\theta G_\eta(x)=0,
\qquad
\mathbb E_\theta[F_\lambda(x)G_\eta(x)]
=\frac{\theta}{2}B_h(x)^2.
\]

Set
\[
c_S:=\frac1{64D_d}.
\]
If the amplitude inequalities hold with \(c_S\), then
\[
a,b\le c_S,\qquad
aD_d,bD_d\le\frac1{64},
\qquad
ab\le\frac1{4096}.
\]
Thus the raw propensity and both raw arm probabilities lie in \([0,1]\), at every covariate, so their Bernoulli atom formulas apply directly.

For a labeled record, define the treatment and observed-response spins by
\[
s^A:=2A-1,\qquad
Y:=AY(1)+(1-A)Y(0),\qquad
s^Y:=2Y-1.
\]
Relative to the uniform probability measure on the four spin pairs, its likelihood is
\[
L_{\theta,\lambda,\eta}(x;s^A,s^Y)
:=
(1+2s^AaF_\lambda(x))
(1+2s^YbG_\eta(x)+s^As^Y\tau_\theta(x)).
\]
This follows by multiplying the two Bernoulli atom probabilities, each divided by its fair atom probability. Expanding the product gives
\[
\begin{aligned}
\mathbb E_\theta L_{\theta,\lambda,\eta}(x;s^A,s^Y)
={}&1+s^As^Y\tau_\theta(x)\\
&+\bigl(2s^Aa+2(s^A)^2s^Ya\tau_\theta(x)\bigr)
  \mathbb E_\theta F_\lambda(x)\\
&+2s^Yb\,\mathbb E_\theta G_\eta(x)
+4s^As^Yab\,\mathbb E_\theta[F_\lambda(x)G_\eta(x)]\\
={}&1+s^As^Y
\bigl(-2\theta abB_h(x)^2+2\theta abB_h(x)^2\bigr)
=1.
\end{aligned}
\]
Hence every atom of the singleton labeled mixture has probability \(1/4\). This proves its uniform conditional law for every hypothesis and every cube covariate.
% lean: CausalSmith.Stat.TwosamplePointcateAnnotationFrontier.singleton_cancellation

\item \textit{Equality on uninformative components.}
For integers \(n\ge2\) and \(m\ge0\), write
\[
N:=n+m,\qquad
\mathbf x:=(x_1,\ldots,x_N)\in([0,1]^d)^N.
\]
The first \(n\) indices are labeled. A spin assignment has domain
\[
\mathbf s\in
\bigl(\{-1,+1\}^2\bigr)^n
\times\{-1,+1\}^m.
\]
Define the record likelihoods by
\[
R_{\theta,\lambda,\eta,i}(\mathbf x,\mathbf s)
:=
\begin{cases}
L_{\theta,\lambda,\eta}(x_i;s_i^A,s_i^Y),&i\le n,\\
1+2s_i^AaF_\lambda(x_i),&i>n,
\end{cases}
\]
and, for every subset \(\mathcal V\subseteq\{1,\ldots,N\}\), define
\[
g_{\theta,\mathcal V}(\mathbf x,\mathbf s)
:=
\mathbb E_\theta
\prod_{i\in\mathcal V}
R_{\theta,\lambda,\eta,i}(\mathbf x,\mathbf s).
\]
Empty products equal one, so \(g_{\theta,\varnothing}=1\).

The sign flip
\[
\mathcal F_\eta(\lambda,\eta):=(\lambda,-\eta)
\]
is a bijection, and its weights satisfy
\[
w_{+1}(\lambda,\eta)=w_{-1}(\lambda,-\eta).
\]
Every auxiliary likelihood is unchanged by this flip, because it depends on \(\lambda\) alone. Reindexing the finite sum consequently gives
\[
g_{+1,\mathcal V}(\mathbf x,\mathbf s)
=
g_{-1,\mathcal V}(\mathbf x,\mathbf s)
\qquad
\text{if }\mathcal V\subseteq\{n+1,\ldots,N\}.
\]
If \(\mathcal V=\{i\}\) with \(i\le n\), singleton cancellation instead gives
\[
g_{+1,\{i\}}(\mathbf x,\mathbf s)
=
g_{-1,\{i\}}(\mathbf x,\mathbf s)=1.
\]
Thus we may take
\[
c_U:=c_S.
\]
Under its amplitude conditions, every auxiliary-only component and every singleton labeled component has equal densities under the two hypotheses. Equality for all spin assignments gives equality of their full conditional laws.
% lean: CausalSmith.Stat.TwosamplePointcateAnnotationFrontier.uninformative_component_equality

\item \textit{The component density difference.}
Define
\[
c_\Delta:=\frac1{64D_d},
\qquad
C_\Delta:=8D_d^2+6.
\]
Assume the amplitude inequalities with \(c_\Delta\). Fix \(\mathbf x,\mathbf s\) and a subset \(\mathcal V\), and define its size by
\[
p_{\mathrm{cmp}}:=|\mathcal V|.
\]
For every record, introduce the contrast-deleted factors
\[
f_{\lambda,i}:=1+2s_i^AaF_\lambda(x_i),
\qquad
g_{\eta,i}:=
\begin{cases}
1+2s_i^YbG_\eta(x_i),&i\le n,\\
1,&i>n.
\end{cases}
\]
The bounds from the previous steps give
\[
|f_{\lambda,i}-1|\le2aD_d,\qquad
|g_{\eta,i}-1|\le2bD_d,
\qquad
|f_{\lambda,i}|,|g_{\eta,i}|\le\frac98.
\]
For labeled records,
\[
R_{\theta,\lambda,\eta,i}
=f_{\lambda,i}
\bigl(g_{\eta,i}+s_i^As_i^Y\tau_\theta(x_i)\bigr),
\qquad
|\tau_\theta(x_i)|\le2ab.
\]
For auxiliary records, \(R_{\theta,\lambda,\eta,i}=f_{\lambda,i}\). Since \(ab\le1/4096\), both cases satisfy
\[
|R_{\theta,\lambda,\eta,i}|
\le\frac32,
\qquad
|R_{\theta,\lambda,\eta,i}-f_{\lambda,i}g_{\eta,i}|
\le3ab,
\qquad
|f_{\lambda,i}g_{\eta,i}|\le\frac32.
\]

Define
\[
\mathcal A_{\mathcal V}(\lambda)
:=\prod_{i\in\mathcal V}f_{\lambda,i},
\qquad
\mathcal G_{\mathcal V}(\eta)
:=\prod_{i\in\mathcal V}g_{\eta,i},
\]
\[
U_{\mathcal V}
:=p_{\mathrm{cmp}}(2aD_d)(9/8)^{p_{\mathrm{cmp}}},
\qquad
W_{\mathcal V}
:=p_{\mathrm{cmp}}(2bD_d)(9/8)^{p_{\mathrm{cmp}}},
\]
\[
D_{\mathcal V}
:=p_{\mathrm{cmp}}(3ab)(3/2)^{p_{\mathrm{cmp}}}.
\]
Telescoping a product, using an upper bound \(D_{\mathrm{prod}}\ge1\) on both sets of factors, gives
\[
\left|
\prod_{i\in\mathcal V}f_{\mathrm{tel},i}
-\prod_{i\in\mathcal V}g_{\mathrm{tel},i}
\right|
\le
D_{\mathrm{prod}}^{|\mathcal V|}
\sum_{i\in\mathcal V}
|f_{\mathrm{tel},i}-g_{\mathrm{tel},i}|.
\]
Here the telescoping sum replaces one factor at a time; each remaining product is bounded by \(D_{\mathrm{prod}}^{|\mathcal V|-1}\), which is at most the displayed power. For the empty set both sides are zero. Applying this inequality gives
\[
|\mathcal A_{\mathcal V}-1|\le U_{\mathcal V},
\qquad
|\mathcal G_{\mathcal V}-1|\le W_{\mathcal V},
\]
\[
\left|
\prod_{i\in\mathcal V}R_{\theta,\lambda,\eta,i}
-\mathcal A_{\mathcal V}\mathcal G_{\mathcal V}
\right|
\le D_{\mathcal V}.
\]

Flipping \(\eta\) preserves every function of \(\lambda\) and changes the positive weights to the negative weights. Flipping \(\lambda\) similarly preserves every function of \(\eta\). Consequently
\[
\mathbb E_{+1}\mathcal A_{\mathcal V}
=\mathbb E_{-1}\mathcal A_{\mathcal V},
\qquad
\mathbb E_{+1}\mathcal G_{\mathcal V}
=\mathbb E_{-1}\mathcal G_{\mathcal V}.
\]
Expanding the centered product and using normalization therefore yields
\[
\begin{aligned}
&\mathbb E_{+1}[\mathcal A_{\mathcal V}\mathcal G_{\mathcal V}]
-\mathbb E_{-1}[\mathcal A_{\mathcal V}\mathcal G_{\mathcal V}]\\
&\qquad=
\mathbb E_{+1}[(\mathcal A_{\mathcal V}-1)(\mathcal G_{\mathcal V}-1)]
-\mathbb E_{-1}[(\mathcal A_{\mathcal V}-1)(\mathcal G_{\mathcal V}-1)].
\end{aligned}
\]
Each expectation on the right has absolute value at most
\(U_{\mathcal V}W_{\mathcal V}\). Adding the two contrast-deletion errors gives
\[
|g_{+1,\mathcal V}-g_{-1,\mathcal V}|
\le D_{\mathcal V}+2U_{\mathcal V}W_{\mathcal V}+D_{\mathcal V}.
\]
Finally,
\[
(9/8)^{p_{\mathrm{cmp}}}(9/8)^{p_{\mathrm{cmp}}}
\le(3/2)^{p_{\mathrm{cmp}}},
\qquad
p_{\mathrm{cmp}}\le p_{\mathrm{cmp}}^2
\quad(p_{\mathrm{cmp}}\in\mathbb N).
\]
Substitution proves
\[
|g_{+1,\mathcal V}(\mathbf x,\mathbf s)
-g_{-1,\mathcal V}(\mathbf x,\mathbf s)|
\le
C_\Delta(3/2)^{p_{\mathrm{cmp}}}
p_{\mathrm{cmp}}^2ab.
\]
For \(p_{\mathrm{cmp}}=0\), both densities equal one and the bound is zero.
% lean: CausalSmith.Stat.TwosamplePointcateAnnotationFrontier.component_information_bound

\item \textit{A positive floor and component Hellinger distance.}
Define
\[
c_H:=\min\left\{c_\Delta,\frac1{64D_d}\right\},
\qquad
K_{\mathrm{info}}:=C_\Delta^2.
\]
Under the amplitude inequalities with \(c_H\),
\[
aD_d\le\frac1{16},
\qquad
bD_d+ab\le\frac1{16}.
\]
Thus the treatment factor of a labeled likelihood is at least \(7/8\), and its outcome factor is also at least \(7/8\). Their product is at least \(1/2\). An auxiliary likelihood is at least \(7/8\), hence at least \(1/2\). Averaging the resulting product floor with nonnegative normalized weights gives
\[
g_{\theta,\mathcal V}(\mathbf x,\mathbf s)
\ge(1/2)^{p_{\mathrm{cmp}}}.
\]

Define the common probability reference on the record spins by
\[
\mu_{\mathrm{spin}}
:=
\bigotimes_{i=1}^{n}\operatorname{Unif}(\{-1,+1\}^2)
\otimes
\bigotimes_{i=n+1}^{N}\operatorname{Unif}(\{-1,+1\}),
\]
and define
\[
H_{\mathcal V}^2(\mathbf x)
:=
\int
\left(
\sqrt{g_{+1,\mathcal V}(\mathbf x,\mathbf s)}
-\sqrt{g_{-1,\mathcal V}(\mathbf x,\mathbf s)}
\right)^2
\,d\mu_{\mathrm{spin}}(\mathbf s).
\]
For nonnegative numbers bounded below by \(t_{\mathrm{floor}}>0\), the identity
\[
(f_{\mathrm{den}}-g_{\mathrm{den}})^2
=
(\sqrt{f_{\mathrm{den}}}-\sqrt{g_{\mathrm{den}}})^2
(\sqrt{f_{\mathrm{den}}}+\sqrt{g_{\mathrm{den}}})^2
\]
and the denominator bound
\[
(\sqrt{f_{\mathrm{den}}}+\sqrt{g_{\mathrm{den}}})^2
\ge4t_{\mathrm{floor}}
\]
give
\[
(\sqrt{f_{\mathrm{den}}}-\sqrt{g_{\mathrm{den}}})^2
\le
\frac{(f_{\mathrm{den}}-g_{\mathrm{den}})^2}{4t_{\mathrm{floor}}}.
\]
Apply this with the preceding component floor and density difference, and integrate against the probability reference. We obtain
\[
\begin{aligned}
H_{\mathcal V}^2(\mathbf x)
&\le
\frac{
\bigl(C_\Delta(3/2)^{p_{\mathrm{cmp}}}
p_{\mathrm{cmp}}^2ab\bigr)^2
}{
4(1/2)^{p_{\mathrm{cmp}}}
}\\
&=
\frac{C_\Delta^2}{4}(9/2)^{p_{\mathrm{cmp}}}
p_{\mathrm{cmp}}^4a^2b^2\\
&\le
K_{\mathrm{info}}(9/2)^{p_{\mathrm{cmp}}}
p_{\mathrm{cmp}}^4a^2b^2.
\end{aligned}
\]
% lean: CausalSmith.Stat.TwosamplePointcateAnnotationFrontier.marked_component_hellinger_bound

\item \textit{Conditional factorization and the count bound.}
For fixed covariates, define the augmented graph
\[
\mathscr G_{\mathbf x}
:=
\left(
\{1,\ldots,N\},
\left\{\{i,j\}:
i\ne j,\ x_i,x_j\in C_h,\
\|x_i-x_j\|_\infty\le2\delta
\right\}
\right),
\]
and let
\[
\mathscr K_{\mathbf x}
:=
\{\mathcal V:\mathcal V
\text{ is the vertex set of a connected component of }
\mathscr G_{\mathbf x}\}.
\]
The cube \(C_h\) is that of \cref{obj:synth_3}. Records outside \(C_h\) are isolated in this augmented graph and have likelihood one.

If \(\phi_{\mathbf z}(x_i)\ne0\), then
\[
x_i\in C_h,
\qquad
\left|\frac{x_{i,j_{\mathrm{coord}}}-x_{0,j_{\mathrm{coord}}}}{\delta}-z_{j_{\mathrm{coord}}}\right|<1
\quad(j_{\mathrm{coord}}\in\{1,\ldots,d\}).
\]
Two records sharing a nonzero frame entry therefore satisfy
\[
\|x_i-x_j\|_\infty\le2\delta,
\]
and belong to the same component. Assign each frame index active at any record to this unique component; inactive indices contribute a factor one. The product prior is independent across indices, and each component likelihood depends solely on its assigned indices. Regrouping the finite sign sums by these disjoint groups consequently proves
\[
g_\theta(\mathbf x,\mathbf s)
:=
\mathbb E_\theta
\prod_{i=1}^{N}R_{\theta,\lambda,\eta,i}(\mathbf x,\mathbf s)
=
\prod_{\mathcal V\in\mathscr K_{\mathbf x}}
g_{\theta,\mathcal V}(\mathbf x,\mathbf s).
\]
Each component density depends solely on its component spins. Integrating the fixed-sign record product over those fair spins gives one, because every record likelihood is normalized. Averaging signs therefore gives
\[
g_{\theta,\mathcal V}\ge0,
\qquad
\int g_{\theta,\mathcal V}\,d\mu_{\mathrm{spin}}=1.
\]

Define the component affinities by
\[
\mathcal A_{\mathrm{aff},\mathcal V}(\mathbf x)
:=
\int\sqrt{g_{+1,\mathcal V}g_{-1,\mathcal V}}\,
d\mu_{\mathrm{spin}}.
\]
Normalization and \(2\sqrt{f_{\mathrm{den}}g_{\mathrm{den}}}
\le f_{\mathrm{den}}+g_{\mathrm{den}}\) imply
\[
0\le\mathcal A_{\mathrm{aff},\mathcal V}\le1,
\qquad
H_{\mathcal V}^2
=2(1-\mathcal A_{\mathrm{aff},\mathcal V}).
\]
Integration over disjoint spin blocks multiplies the affinities. The elementary finite-product inequality
\[
1-\prod_{\mathcal V\in\mathscr K_{\mathbf x}}
\mathcal A_{\mathrm{aff},\mathcal V}
\le
\sum_{\mathcal V\in\mathscr K_{\mathbf x}}
(1-\mathcal A_{\mathrm{aff},\mathcal V})
\]
follows by induction from the fact that all the affinities lie in \([0,1]\). Hence
\[
\int(\sqrt{g_{+1}}-\sqrt{g_{-1}})^2\,d\mu_{\mathrm{spin}}
\le
\sum_{\mathcal V\in\mathscr K_{\mathbf x}}
H_{\mathcal V}^2(\mathbf x).
\]

Set
\[
c_I:=\min\{c_H,c_U\}.
\]
Its amplitude conditions permit both the component bound and the uninformative-component equalities. A component containing a labeled record and having size less than two has size one, so its distance is zero by singleton cancellation. A component containing no labeled record is auxiliary-only, so its distance is also zero. Define the informative component counts, for every integer \(p_{\mathrm{cmp}}\ge0\), by
\[
C_{p_{\mathrm{cmp}}}^{\mathrm L}(\mathbf x)
:=
\#\left\{
\mathcal V\in\mathscr K_{\mathbf x}:
|\mathcal V|=p_{\mathrm{cmp}},\
\mathcal V\cap\{1,\ldots,n\}\ne\varnothing
\right\}.
\]
Regrouping the remaining component sum by size gives
\[
\int(\sqrt{g_{+1}}-\sqrt{g_{-1}})^2\,d\mu_{\mathrm{spin}}
\le
K_{\mathrm{info}}a^2b^2
\sum_{p_{\mathrm{cmp}}=2}^{N}
(9/2)^{p_{\mathrm{cmp}}}
p_{\mathrm{cmp}}^4C_{p_{\mathrm{cmp}}}^{\mathrm L}(\mathbf x).
\]
The graph \(\mathscr G_{\mathbf x}\) is exactly the graph in the statement: it retains all \(N\) vertices, including the exterior singleton components, whose density factors equal one.
% lean: CausalSmith.Stat.TwosamplePointcateAnnotationFrontier.marked_conditional_hellinger_count_bound

\item \textit{The labeled-root tree series.}
Define the common covariate probability measure by
\[
\nu_{\mathrm{cov}}
:=
\operatorname{Unif}([0,1]^d)^{\otimes N}.
\]
For \(p_{\mathrm{cmp}}\ge2\), a component counted by
\(C_{p_{\mathrm{cmp}}}^{\mathrm L}\) has a labeled root and a spanning tree on its vertex set. For a fixed root, there are
\(\binom{N-1}{p_{\mathrm{cmp}}-1}\) choices of the other vertices. Orient a spanning tree toward the root. Each nonroot vertex chooses a parent among the \(p_{\mathrm{cmp}}\) selected vertices, so the number of valid rooted trees is at most
\(p_{\mathrm{cmp}}^{p_{\mathrm{cmp}}-1}\).

For a fixed valid tree, integrate its leaf covariates successively. Each child is constrained to a maximum-norm box of side \(4\delta\) around its parent; intersection with the unit cube has uniform mass at most \((4\delta)^d\). The root lies in \(C_h\), whose uniform mass is \(h^d\). Thus the probability of this tree event is at most
\[
h^d\bigl((4\delta)^d\bigr)^{p_{\mathrm{cmp}}-1}.
\]
Dropping the remaining component conditions enlarges the event. Summing these tree witnesses over the \(n\) possible labeled roots gives
\[
\int C_{p_{\mathrm{cmp}}}^{\mathrm L}(\mathbf x)\,
d\nu_{\mathrm{cov}}(\mathbf x)
\le
nh^d
\binom{N-1}{p_{\mathrm{cmp}}-1}
p_{\mathrm{cmp}}^{p_{\mathrm{cmp}}-1}
\bigl((4\delta)^d\bigr)^{p_{\mathrm{cmp}}-1}.
\]
For \(p_{\mathrm{cmp}}>N\), the count and the binomial coefficient are both zero.

The binomial bound and the exponential series imply
\[
\binom{N-1}{p_{\mathrm{cmp}}-1}
\le\frac{N^{p_{\mathrm{cmp}}-1}}{(p_{\mathrm{cmp}}-1)!},
\qquad
\frac{p_{\mathrm{cmp}}^{p_{\mathrm{cmp}}-1}}
{(p_{\mathrm{cmp}}-1)!}
\le \exp(p_{\mathrm{cmp}}).
\]
For the second inequality, its left side is one nonnegative summand of the exponential series at \(p_{\mathrm{cmp}}\). Consequently
\[
\binom{N-1}{p_{\mathrm{cmp}}-1}
p_{\mathrm{cmp}}^{p_{\mathrm{cmp}}-1}
\le N^{p_{\mathrm{cmp}}-1}\exp(p_{\mathrm{cmp}}).
\]

Define
\[
\kappa_{\mathrm{tree}}:=\frac92,
\qquad
S_{\mathrm{tree}}
:=
\sum_{j=0}^{\infty}(j+2)^4\,2^{-j},
\qquad
E_{\mathrm{occ}}
:=
\kappa_{\mathrm{tree}}\exp(1)4^d,
\]
\[
c_0:=\frac1{2E_{\mathrm{occ}}},
\qquad
C_{\mathrm{tree}}
:=
(S_{\mathrm{tree}}+1)
\bigl(\kappa_{\mathrm{tree}}\exp(1)\bigr)^2\,4^d.
\]
The series defining \(S_{\mathrm{tree}}\) converges because it is a shifted polynomial-geometric series. These definitions give \(c_0>0\) and \(C_{\mathrm{tree}}>0\).

Under \(N\delta^d\le c_0\), define
\[
z_{\mathrm{occ}}:=E_{\mathrm{occ}}N\delta^d,
\qquad
0\le z_{\mathrm{occ}}\le\frac12.
\]
Applying the binomial-tree bound with \(p_{\mathrm{cmp}}=j+2\) gives
\[
\begin{aligned}
&\kappa_{\mathrm{tree}}^{j+2}(j+2)^4
\binom{N-1}{j+1}(j+2)^{j+1}
\bigl((4\delta)^d\bigr)^{j+1}\\
&\qquad\le
\kappa_{\mathrm{tree}}\exp(1)(j+2)^4
z_{\mathrm{occ}}^{j+1}.
\end{aligned}
\]
For every finite truncation,
\[
\sum_{j=0}^{N-1}(j+2)^4z_{\mathrm{occ}}^{j+1}
\le
z_{\mathrm{occ}}\sum_{j=0}^{N-1}(j+2)^4\,2^{-j}
\le z_{\mathrm{occ}}S_{\mathrm{tree}}.
\]
Multiplying the tree-count inequality by its size weight and summing therefore yields
\[
\begin{aligned}
&\sum_{j=0}^{N-1}
\kappa_{\mathrm{tree}}^{j+2}(j+2)^4
\int C_{j+2}^{\mathrm L}\,d\nu_{\mathrm{cov}}\\
&\qquad\le
nh^d\kappa_{\mathrm{tree}}\exp(1)
z_{\mathrm{occ}}S_{\mathrm{tree}}\\
&\qquad\le C_{\mathrm{tree}}nNh^d\delta^d.
\end{aligned}
\]
The final inequality replaces \(S_{\mathrm{tree}}\) by
\(S_{\mathrm{tree}}+1\) and substitutes the definition of
\(z_{\mathrm{occ}}\).
% lean: CausalSmith.Stat.TwosamplePointcateAnnotationFrontier.labeled_component_series_bound

\item \textit{Integration for the original experiment.}
The covariate distribution is \(\nu_{\mathrm{cov}}\) for every sign realization and hypothesis. Conditional on these covariates, the original records have density \(g_\theta\) against \(\mu_{\mathrm{spin}}\). Thus, after grouping covariates and spins, the mixture laws have densities \(g_{+1}\) and \(g_{-1}\) against
\(\nu_{\mathrm{cov}}\otimes\mu_{\mathrm{spin}}\).

For two densities against a common dominating measure, substitution into \cref{obj:def:hellinger} gives the integral of their squared square-root difference. Applying this identity and then integrating over covariates yields
\[
H^2(M_{+1},M_{-1})
=
\int
\left[
\int(\sqrt{g_{+1}(\mathbf x,\mathbf s)}
-\sqrt{g_{-1}(\mathbf x,\mathbf s)})^2
\,d\mu_{\mathrm{spin}}(\mathbf s)
\right]d\nu_{\mathrm{cov}}(\mathbf x).
\]
The finite graph counts are measurable, and all functions in the finite weighted count sum are integrable. Integrating the conditional bound therefore gives
\[
H^2(M_{+1},M_{-1})
\le
K_{\mathrm{info}}a^2b^2
\sum_{p_{\mathrm{cmp}}=2}^{N}
(9/2)^{p_{\mathrm{cmp}}}p_{\mathrm{cmp}}^4
\int C_{p_{\mathrm{cmp}}}^{\mathrm L}\,d\nu_{\mathrm{cov}}.
\]
Every term is nonnegative, so extending the size sum to the shifted range used above gives
\[
\begin{aligned}
H^2(M_{+1},M_{-1})
&\le
K_{\mathrm{info}}a^2b^2
\sum_{j=0}^{N-1}
(9/2)^{j+2}(j+2)^4
\int C_{j+2}^{\mathrm L}\,d\nu_{\mathrm{cov}}\\
&\le
K_{\mathrm{info}}C_{\mathrm{tree}}
nNh^d\delta^da^2b^2.
\end{aligned}
\]
Define
\[
C:=K_{\mathrm{info}}C_{\mathrm{tree}}>0.
\]
This proves the required information estimate under the amplitude conditions with \(c_I\) and the occupancy condition with \(c_0\).
% lean: CausalSmith.Stat.TwosamplePointcateAnnotationFrontier.marked_information_bound

\item \textit{Common calibration and the remaining equalities.}
Choose
\[
c:=\min\{c_M,\min\{c_S,\min\{c_U,c_I\}\}\}.
\]
All four constants are positive and at most one, so
\[
0<c\le1,
\qquad
c\le c_M,\quad c\le c_S,\quad c\le c_U,\quad c\le c_I.
\]
Fix a tuple satisfying the theorem's hypotheses with this \(c\). Since \(\delta>0\) and \(h>0\), multiplication by their real powers preserves these comparisons. For each
\(c'\in\{c_M,c_S,c_U,c_I\}\),
\[
a\le c\delta^\alpha\le c'\delta^\alpha,\qquad
b\le c\delta^\beta\le c'\delta^\beta,\qquad
ab\le ch^\gamma\le c'h^\gamma.
\]
The membership, singleton, uninformative-component, and information estimates established above therefore apply simultaneously. Weight normalization makes the finite pushforward priors probability measures, and the pointwise membership conclusion supports them on \(\mathcal M\).

For the propensity distribution, the sign flip preserves the covariate-treatment marginal:
\[
(P_{+1,\lambda,\eta})_{XA}
=
(P_{-1,\lambda,-\eta})_{XA}.
\]
The designated propensity witnesses of these admitted laws are continuous and are versions of the same conditional treatment probability. Uniqueness of conditional margins followed by continuity under the full-support uniform design gives
\[
e_{P_{+1,\lambda,\eta}}(x)
=
e_{P_{-1,\lambda,-\eta}}(x)
\qquad(x\in[0,1]^d).
\]
Reindexing the finite pushforwards using
\(w_{+1}(\lambda,\eta)=w_{-1}(\lambda,-\eta)\) proves equality of the distributions of the entire propensity functions restricted to the cube.

At the target, \(B_h(x_0)=1\), and hence
\[
\tau_\theta(x_0)=-2\theta ab,
\qquad
\tau_{-1}(x_0)-\tau_{+1}(x_0)=4ab.
\]
Because \(a,b>0\) and \(c\le1\le4\),
\[
cab=c(ab)\le4(ab)=4ab.
\]

Finally, for every integer \(k_{\mathrm{aux}}\ge0\), reindexing by the same sign flip gives
\[
\begin{aligned}
\int P_{XA}^{\otimes k_{\mathrm{aux}}}\,d\varpi_{+1}(P)
&=
\sum_{\lambda,\eta}
w_{+1}(\lambda,\eta)
(P_{+1,\lambda,\eta})_{XA}^{\otimes k_{\mathrm{aux}}}\\
&=
\sum_{\lambda,\eta}
w_{-1}(\lambda,\eta)
(P_{-1,\lambda,\eta})_{XA}^{\otimes k_{\mathrm{aux}}}\\
&=
\int P_{XA}^{\otimes k_{\mathrm{aux}}}\,d\varpi_{-1}(P).
\end{aligned}
\]
For any fixed auxiliary covariates, their conditional mixture kernels are
\[
\sum_{\lambda,\eta}
w_\theta(\lambda,\eta)
\bigotimes_{i=1}^{k_{\mathrm{aux}}}
\operatorname{Bernoulli}(e_\lambda(x_i)).
\]
The factors are unchanged by the sign flip, so these kernels also agree for every covariate collection. When \(k_{\mathrm{aux}}=0\), the products are the unit probability measure on the empty tuple, and normalization gives the same equality.

The singleton calculation supplies the stated uniform labeled conditional laws. The conditional factorization and component equalities apply to every cube covariate collection and every spin assignment. Together with the calibrated information bound, these establish all assertions.
% lean: CausalSmith.Stat.TwosamplePointcateAnnotationFrontier.marked_component_certificate
\end{enumerate}
\end{proof}
